% ==========================================================================
% Ouvidoria da Mulher no WhatsApp (piloto TJPA / EscutAI)
% Planejamento de atendimento, Chatwoot e implementação. NADA FOI IMPLEMENTADO.
% Compilar: pdflatex Ouvidoria_Mulher_Planejamento.tex (duas vezes)
% ==========================================================================
\documentclass[aspectratio=169,10pt]{beamer}

\usepackage[utf8]{inputenc}
\usepackage[T1]{fontenc}
\usepackage[brazil]{babel}
\usepackage[scaled=0.95]{helvet}
\renewcommand{\familydefault}{\sfdefault}
\usepackage{microtype}
\usepackage{graphicx}
\usepackage{booktabs}
\usepackage{tabularx}
\usepackage{array}
\usepackage{fontawesome5}
\usepackage{tikz}
\usetikzlibrary{arrows.meta,positioning,calc,shapes.geometric,fit,backgrounds}
\usepackage{hyperref}
\hypersetup{colorlinks=true,urlcolor=petroleo,linkcolor=petroleo}
\graphicspath{{../TJPA/img/}}

% Identidade visual: mesma paleta do pitch EscutAI + rosa da Cartilha da Ouvidoria da Mulher
\definecolor{escuro}{HTML}{053F35}
\definecolor{petroleo}{HTML}{0B5345}
\definecolor{verde}{HTML}{00A884}
\definecolor{menta}{HTML}{16E0A9}
\definecolor{ambar}{HTML}{FDB833}
\definecolor{areia}{HTML}{FDE3A7}
\definecolor{nevoa}{HTML}{E8F4F0}
\definecolor{tinta}{HTML}{0E2F33}
\definecolor{cinza}{HTML}{5E6B6B}
\definecolor{seta}{HTML}{A8D5C6}
\definecolor{rosa}{HTML}{B03060}
\definecolor{rosaclaro}{HTML}{F8E1EA}
\definecolor{alerta}{HTML}{C0392B}

\setbeamercolor{normal text}{fg=tinta,bg=white}
\setbeamercolor{structure}{fg=petroleo}
\setbeamercolor{frametitle}{fg=tinta}
\setbeamercolor{itemize item}{fg=rosa}
\setbeamercolor{itemize subitem}{fg=ambar}
\setbeamercolor{alerted text}{fg=rosa}
\setbeamerfont{frametitle}{size=\Large,series=\bfseries}
\setbeamertemplate{navigation symbols}{}
\setbeamertemplate{itemize item}{\footnotesize\raise0.1ex\hbox{\textcolor{rosa}{\faCaretRight}}}
\setbeamertemplate{itemize subitem}{\scriptsize\textcolor{ambar}{\faCircle}}
\setbeamersize{text margin left=0.75cm,text margin right=0.75cm}

\setbeamertemplate{frametitle}{%
  \vskip0.30cm
  \begin{minipage}[c]{0.74\paperwidth}\usebeamerfont{frametitle}\usebeamercolor[fg]{frametitle}\insertframetitle\end{minipage}%
  \hfill
  \begin{minipage}[c]{0.15\paperwidth}\raggedleft
    \includegraphics[height=0.55cm]{logo_escutai_verde.png}\hspace{0.1cm}{\bfseries\small\textcolor{tinta}{Escut}\textcolor{verde}{AI}}%
  \end{minipage}\par
  \vskip-0.10cm
}
\setbeamertemplate{footline}{%
  \leavevmode\hbox to\paperwidth{%
    \hspace{0.75cm}{\tiny\textcolor{cinza}{Ouvidoria da Mulher no WhatsApp \textperiodcentered\ planejamento (nada implementado)}}\hfill
    {\tiny\textcolor{cinza}{\insertframenumber}}\hspace{0.75cm}}\vskip0.25cm}

\newcommand{\fundoescuro}{%
  \setbeamercolor{background canvas}{bg=escuro}%
  \setbeamercolor{normal text}{fg=white}%
  \setbeamertemplate{footline}{}%
  \usebeamercolor[fg]{normal text}}

\newcommand{\code}[1]{\texttt{\detokenize{#1}}}
\newcommand{\fonte}[1]{\par\vskip3pt{\tiny\textcolor{cinza}{Fonte: #1}}}

\tikzset{
  caixa/.style={rounded corners=5pt,align=center,inner sep=4pt,font=\scriptsize},
  nova/.style={caixa,fill=rosa,text=white,font=\scriptsize\bfseries},
  atual/.style={caixa,fill=nevoa,draw=verde,text=tinta},
  externa/.style={caixa,fill=areia,text=tinta},
  seta/.style={-{Stealth[length=5pt]},line width=1pt,color=petroleo},
  setanova/.style={-{Stealth[length=5pt]},line width=1.2pt,color=rosa},
}
\renewcommand{\arraystretch}{1.18}

\title{Ouvidoria da Mulher no WhatsApp}
\subtitle{Planejamento de atendimento, Chatwoot e implementação}
\author{Equipe EscutAI}
\date{Piloto TJPA}

\begin{document}

% ==========================================================================
% CAPA
% ==========================================================================
{\fundoescuro
\begin{frame}[plain]
  \vspace{1.4cm}
  {\color{menta}\bfseries\small PILOTO TJPA \textperiodcentered\ ESCUTAI \textperiodcentered\ JÚLIA}\\[0.3cm]
  {\fontsize{28}{32}\selectfont\bfseries Ouvidoria da Mulher\\no WhatsApp}\\[0.35cm]
  {\large Planejamento do atendimento, o que fazer no Chatwoot\\e plano de implementação}\\[0.9cm]
  \begin{tikzpicture}
    \node[fill=ambar,text=tinta,rounded corners=4pt,inner sep=6pt,font=\small\bfseries,align=left]
      {\faExclamationTriangle\ Documento de planejamento. Nenhum código, Chatwoot ou VM foi alterado.\\
       A implementação só começa com a sua autorização.};
  \end{tikzpicture}
\end{frame}}

% ==========================================================================
% 1. RESUMO
% ==========================================================================
\begin{frame}{Resumo para decisão}
\begin{itemize}\setlength\itemsep{3pt}\small
  \item \alert{O quê.} Novo serviço \textbf{Ouvidoria da Mulher} no \textbf{menu principal existente} da Júlia. Ele orienta com base na Cartilha da Ouvidoria da Mulher e faz uma \textbf{triagem estruturada} para a equipe humana.
  \item \alert{A triagem não resolve nem mede risco.} Coleta 8 dados mínimos, todos com a opção \emph{prefiro não responder}, e entrega ao atendente numa \textbf{nota privada} do Chatwoot, mais atributos e rótulos.
  \item \alert{Chatwoot.} Criar um \textbf{novo time}, \code{ouvidoria da mulher}, na caixa de entrada WhatsApp atual. \textbf{Novo canal} só se vocês exigirem isolamento de acesso ou número próprio (pág.~\ref{pg:time}).
  \item \alert{Segurança primeiro.} 190 e 180 sempre visíveis, a pergunta de perigo vem antes das demais, nenhuma promessa de prazo nem de sigilo absoluto, sem CPF, endereço ou dados do agressor.
  \item \alert{Entrega.} Atrás de uma chave de configuração \textbf{desligada por padrão}. O deploy leva só código; ligar é decisão sua.
  \item \alert{O que preciso de você.} 9 decisões na pág.~\ref{pg:decisoes}. Os contatos de emergência precisam de confirmação em fonte oficial antes de publicar.
\end{itemize}
\end{frame}

% ==========================================================================
% A. BASE DE EVIDÊNCIAS
% ==========================================================================
\begin{frame}{Cartilha: o que vira conteúdo do canal}
\scriptsize
\begin{tabularx}{\linewidth}{@{}>{\bfseries}p{2.9cm}p{1.2cm}>{\raggedright\arraybackslash}X>{\raggedright\arraybackslash}X@{}}
\toprule
Seção & Págs. & O que a Cartilha traz & Uso no canal \\
\midrule
O que é e o que faz & 4, 19 & Espaço seguro de escuta, acolhimento e orientação. Recebe denúncias, atua em morosidade processual, escuta, orienta e encaminha. & Texto de acolhimento e \emph{o que a Ouvidoria faz}. \\
Formas de violência & 6--9 & Física, sexual, psicológica, patrimonial, moral, on-line, perseguição (\emph{stalking}, art.~147-A do CP), tráfico, feminicídio. & Opções de \emph{assunto} da triagem e definições curtas. \\
Assédio e importunação & 11--13 & Moral; sexual (art.~216-A do CP); importunação sexual (art.~215-A do CP). & Opção \emph{assédio/importunação} e explicação da diferença. \\
Leis importantes & 15--17 & Lei Maria da Penha (11.340/06, art.~5º), Lei 14.550/2023, Lei do Feminicídio (13.104/2015). & Orientações e opções de \emph{ambiente}. \\
Como denunciar & 21 & Seis passos; Ligue 180 (24 h, gratuito, confidencial); denúncias tratadas com \emph{sigilo profissional}. & Orientações, \alert{com ressalva}: os passos 2 e 5 (RH, sindicato, Delegacia Regional do Trabalho) valem só para assédio no trabalho. \\
\bottomrule
\end{tabularx}
\vskip4pt
{\scriptsize\textcolor{cinza}{A Cartilha \textbf{não} traz o 190 nem os contatos da Ouvidoria do TJPA: eles vêm da Carta de Serviços e de fontes externas (próximas páginas). O PDF parece ser só imagem; a ingestão exigirá OCR.}}
\end{frame}

\begin{frame}{Ouvidoria da Mulher do TJPA: o que é oficial}
\begin{columns}[T]
\begin{column}{0.52\textwidth}
\begin{tikzpicture}
  \node[fill=rosaclaro,rounded corners=6pt,inner sep=8pt,text width=0.92\textwidth,font=\small,align=left]{%
    \textbf{Competência}\\[2pt]
    Receber e encaminhar às autoridades competentes demandas sobre \textbf{procedimentos judiciais} de violência contra a mulher. Recebe também informações, sugestões, reclamações, denúncias, críticas e elogios sobre a \textbf{tramitação} desses procedimentos.\\[6pt]
    \textbf{Contatos}\\[2pt]
    \faPhone\ (91) 3289-7182\\
    \faEnvelope\ ouvidoria.mulher@tjpa.jus.br\\
    \faClock\ \textbf{8h às 14h}};
\end{tikzpicture}
\fonte{Carta de Serviços do TJPA, ficha \emph{Ouvidoria da Mulher} (PDF oficial lido). A data da edição não aparece: confirmar os dados antes de publicar.}
\end{column}
\begin{column}{0.46\textwidth}
\textbf{O que isso muda no desenho}
\begin{itemize}\setlength\itemsep{4pt}\small
  \item \alert{Não é canal de emergência.} Atende 8h--14h: o bot não pode sugerir resposta imediata e mostra 190 e 180 sempre.
  \item \alert{O escopo é processual.} A triagem precisa captar \emph{existe processo ou medida protetiva?} e o número, quando houver.
  \item \alert{Fora do escopo} (ex.: crime em andamento, sem processo): orientar Ligue 180, delegacia (DEAM) e 190, como a Cartilha (p.~21).
\end{itemize}
\end{column}
\end{columns}
\end{frame}

\begin{frame}{Evidências externas e efeito no desenho}
\fontsize{6.2}{7.4}\selectfont
\begin{tabularx}{\linewidth}{@{}>{\raggedright\arraybackslash}p{2.6cm}>{\raggedright\arraybackslash}X>{\raggedright\arraybackslash}X>{\raggedright\arraybackslash}p{2.0cm}@{}}
\toprule
Fonte & O que afirma & Decisão de desenho & Verificação \\
\midrule
OMS, manual clínico (\emph{LIVES}) & Primeira linha: ouvir, perguntar sobre necessidades sem pressionar, validar, reforçar a segurança, apoiar e conectar a serviços; respeitar as escolhas. & Acolhimento antes de perguntas; cada pergunta com \emph{prefiro não responder}; perigo e segurança primeiro; nada de interrogatório. & Resumo da página da OMS lido; PDF completo não aberto. \\
FONAR (Res. Conj. CNJ/CNMP 5/2020; Lei 14.149/2021) & Formulário de risco com 17 dimensões; Parte~I aplicada por Polícia Civil, MP ou Judiciário, e pela própria vítima \emph{só na ausência} de profissional; Parte~II por psicólogo, assistente social ou médico. & O bot \textbf{não aplica o FONAR nem calcula risco}. Coleta só 2 sinais de segurança e encaminha. & Guia do TJPR lido (PDF). Texto atual: ver Portaria Conj.~6/2025. \\
\emph{Danger Assessment} (Campbell et al., 2009) & Instrumento para estimar letalidade; uso por profissionais treinados; guia, não escore exato. Há estudo discordante (Storey, 2014). & Não criar escore próprio no WhatsApp. & Resumo de busca. \\
Lei 14.550/2023 & Medida protetiva a partir da denúncia à polícia ou de alegações por escrito à Justiça; sem prazo fixo, vale enquanto durar o risco. & Pergunta \emph{existe medida protetiva?}; a Júlia não promete nem prevê decisão. & Cartilha p.~16 e fontes secundárias. \\
Ligue 180 e 190 & 180: 24 h, gratuito, anônimo, com WhatsApp; emergência: 190. & Cartão fixo com 190 e 180 em todo o fluxo. & \alert{Número do WhatsApp do 180 vem de fonte secundária: confirmar em gov.br.} \\
NNEDV (segurança digital) & 97\% dos programas relatam abuso por tecnologia, como monitoramento e ameaças. & \emph{Contato seguro?}, saída rápida, aviso de apagar a conversa, sem texto sensível em notificações. & Levantamento da NNEDV (resumo). \\
\bottomrule
\end{tabularx}
\vskip3pt
{\tiny\textcolor{cinza}{Não utilizadas por falta de verificação nesta pesquisa: Leis 14.188/2021 e 14.321/2022. Resolução CNJ 254/2018 (política judiciária) e 492/2023 (perspectiva de gênero) foram achadas só de forma indireta.}}
\end{frame}

\begin{frame}{Princípios que guiam o desenho}
\begin{columns}[T]
\begin{column}{0.48\textwidth}
\begin{itemize}\setlength\itemsep{5pt}\small
  \item \alert{Segurança antes de dados.} A primeira pergunta é sobre perigo agora; 190 e 180 aparecem desde o primeiro acolhimento.
  \item \alert{Triagem não é barreira.} Em qualquer passo, \emph{Falar com atendente} encaminha com o que já foi coletado.
  \item \alert{Ninguém é forçada a responder.} Toda pergunta tem \emph{prefiro não responder}.
  \item \alert{Contar uma vez só.} O que foi dito chega ao atendente na nota; a pessoa não repete a história.
\end{itemize}
\end{column}
\begin{column}{0.48\textwidth}
\begin{itemize}\setlength\itemsep{5pt}\small
  \item \alert{Mínimo necessário.} Sem CPF, endereço, nome ou dados do agressor, nem detalhes da violência; o relato livre é opcional.
  \item \alert{Sem promessas.} Nada de prazo, resultado, prioridade ou sigilo absoluto (regra da persona); a Cartilha fala em \emph{sigilo profissional}.
  \item \alert{Honestidade.} Encaminhar só é dito depois da confirmação técnica do Chatwoot (regra já existente, CHW-10). Processo informado é \emph{informado}, não \emph{localizado}.
  \item \alert{Sem invenção.} Conteúdo vem da Cartilha, da Carta de Serviços ou de fonte oficial citada; nada de fluxo sem base.
\end{itemize}
\end{column}
\end{columns}
\end{frame}

% ==========================================================================
% B. ACOPLAMENTO À SOLUÇÃO ATUAL
% ==========================================================================
\begin{frame}{Como se acopla à solução atual}
\centering
\begin{tikzpicture}[every node/.style={font=\scriptsize}]
  \node[externa] (wa) at (0.9,0) {Cidadã\\WhatsApp};
  \node[atual] (cw) at (3.9,0) {Chatwoot\\caixa WhatsApp};
  \node[atual] (bot) at (7.0,0) {Agent Bot\\webhook};
  \node[atual,text width=3.9cm] (api) at (11.6,0) {Backend FastAPI\\\code{handler.py} $\to$ \code{flow.py}\\estado em SQLite};
  \node[atual,text width=2.6cm] (rag) at (9.6,-1.9) {RAG institucional\\(domínio aprovado)};
  \node[nova,text width=3.2cm] (mod) at (13.0,-1.9) {Novo módulo\\Ouvidoria da Mulher:\\menu, triagem, nota};
  \node[atual,text width=5.2cm] (team) at (7.0,-3.5) {Times: \code{secretaria}, \code{gabinete}\\\textcolor{rosa}{\textbf{+ ouvidoria da mulher}}};
  \draw[seta] (wa) -- (cw); \draw[seta] (cw) -- (bot); \draw[seta] (bot) -- (api);
  \draw[seta] (api.south) -- ++(0,-0.5) -| (rag.north);
  \draw[setanova] (api.south east) -- ++(0.4,-0.5) -| (mod.north);
  \draw[setanova] (mod.south) |- (team.east);
  \node[font=\tiny,text=rosa,anchor=west] at (10.0,-3.15) {nota privada, rótulos, atributos, \code{assign_team}};
\end{tikzpicture}
\vskip6pt
\begin{itemize}\small\setlength\itemsep{2pt}
  \item[] \textbf{Reaproveita:} máquina de estados do atendimento, protocolo \code{TKT}, idempotência por evento, transferência confirmada por equipe e status, encerramento por inatividade (23 h), cadastro com fonte e data.
  \item[] \textbf{Não cria:} segundo orquestrador, novo canal, nem altera \code{secretaria} e \code{gabinete}.
\end{itemize}
\end{frame}

\begin{frame}{Onde entra no menu existente}
\begin{columns}[T]
\begin{column}{0.45\textwidth}
\textbf{Hoje} (lista do WhatsApp, 5 linhas)
\vskip3pt
\begin{tikzpicture}
  \node[draw=cinza,rounded corners=6pt,fill=nevoa,inner sep=6pt,text width=4.6cm,font=\scriptsize,align=left]{%
  \faUser\ Sou parte/testemunha\\ \faBalanceScale\ Sou advogado(a)\\ \faInfoCircle\ Informações gerais\\ \faHandPaper\ Falar com atendente\\ \faCheck\ Encerrar atendimento};
\end{tikzpicture}
\end{column}
\begin{column}{0.53\textwidth}
\textbf{Proposto} (6 linhas; limite do WhatsApp: 10)
\vskip3pt
\begin{tikzpicture}
  \node[draw=rosa,rounded corners=6pt,fill=white,inner sep=6pt,text width=4.9cm,font=\scriptsize,align=left]{%
  \faUser\ Sou parte/testemunha\\ \faBalanceScale\ Sou advogado(a)\\ \faInfoCircle\ Informações gerais\\
  \colorbox{rosaclaro}{\textcolor{rosa}{\faHeart\ \textbf{Ouvidoria da Mulher}}}\\ \faHandPaper\ Falar com atendente\\ \faCheck\ Encerrar atendimento};
\end{tikzpicture}
\end{column}
\end{columns}
\vskip6pt
\begin{itemize}\small\setlength\itemsep{3pt}
  \item Título da linha com até 24 caracteres: \code{Ouvidoria da Mulher} mais o ícone cabe. O clique volta como texto e é reconhecido como as demais opções.
  \item \textbf{Chave de configuração} desligada por padrão: sem ela, o menu fica idêntico ao de hoje (teste de regressão).
  \item \textbf{Atalho por texto livre:} frases como \emph{meu marido me bate} oferecem a Ouvidoria e mostram 190 e 180; \alert{só oferecem, nunca forçam} (evita falso positivo em perguntas sobre processos).
  \item Visão de futuro: o \emph{Mosaico de Serviços} do TJPA tem 28 serviços; este módulo é o primeiro e serve de molde para os próximos.
\end{itemize}
\end{frame}

\begin{frame}{Mapa do atendimento}
\centering
\begin{tikzpicture}[every node/.style={font=\scriptsize}]
  \node[nova,text width=1.8cm] (a) at (1.1,0) {Acolhimento\\+ 190 e 180};
  \node[nova,text width=1.3cm] (b) at (3.9,0) {Perigo\\agora?};
  \node[nova,text width=2.7cm] (d) at (7.4,0) {Triagem: para quem, assunto, ambiente, processo, contato};
  \node[nova,text width=1.7cm] (e) at (10.6,0) {Relato livre\\(opcional)};
  \node[nova,text width=1.7cm] (f) at (13.3,0) {Resumo e\\consentimento};
  \node[atual,text width=3.0cm] (c) at (5.8,1.9) {Cartão de emergência: 190 e 180; urgência marcada};
  \node[atual,text width=2.3cm] (o) at (1.1,-1.9) {Orientações\\(Cartilha, com página)};
  \node[atual,text width=2.6cm] (x) at (7.4,-1.9) {Contato inseguro: telefone e e-mail, sem encaminhar};
  \node[atual,text width=2.3cm] (g) at (13.3,-1.9) {Transferência para o time \code{ouvidoria da mulher}};
  \draw[seta] (a)--(b); \draw[setanova] (b.north)|-node[above,pos=0.75,font=\tiny]{sim}(c.west);
  \draw[seta] (b)--node[above,font=\tiny]{não}(d); \draw[seta] (c.east)-|(d.north);
  \draw[seta] (d)--(e); \draw[seta] (e)--(f); \draw[setanova] (f)--node[right,font=\tiny]{encaminhar}(g);
  \draw[seta] (a)--(o); \draw[seta] (d)--(x);
\end{tikzpicture}
\vskip6pt
{\scriptsize\textbf{Atalhos em todo o fluxo:} \faHandPaper\ \emph{Falar com atendente} encaminha com o que já existe \textperiodcentered\ \faArrowLeft\ \emph{Retornar ao menu} sai sem enviar \textperiodcentered\ digitar \emph{sair} também sai (saída rápida).}
\end{frame}

\begin{frame}{Orientações: Cartilha com página citada}
\begin{columns}[T]
\begin{column}{0.50\textwidth}
\scriptsize
\begin{tabularx}{\linewidth}{@{}>{\bfseries}p{3.1cm}>{\raggedright\arraybackslash}X@{}}
\toprule
Linha da lista & Fonte \\
\midrule
O que é a Ouvidoria & Cartilha p.~4, 19; Carta de Serviços \\
Formas de violência & p.~6--9 \\
Assédio: qual a diferença & p.~10--13 \\
Lei Maria da Penha & p.~15--16 \\
Lei do Feminicídio & p.~17 \\
Como denunciar & p.~21, com ressalva \\
Telefones úteis & 180 e 190; Ouvidoria (Carta) \\
Outra dúvida & busca restrita à Cartilha \\
\bottomrule
\end{tabularx}
\end{column}
\begin{column}{0.46\textwidth}
\begin{itemize}\small\setlength\itemsep{4pt}
  \item \alert{Texto fixo e aprovado} pela Ouvidoria para os 7 primeiros itens: resposta idêntica e previsível, sem improviso do modelo.
  \item \alert{\emph{Outra dúvida}:} busca só nos trechos da Cartilha aprovados; se não achar, diz que não achou e oferece a Ouvidoria.
  \item Nunca dá orientação jurídica individual ("você tem direito a..."), nem prevê decisão.
  \item A ingestão da Cartilha entra como \code{pending_review}; \textbf{um humano aprova} (regra do projeto).
\end{itemize}
\end{column}
\end{columns}
\end{frame}

% ==========================================================================
% C. TRIAGEM
% ==========================================================================
\begin{frame}{Triagem: 8 passos, todos com saída}
\fontsize{6.2}{7.3}\selectfont
\begin{tabularx}{\linewidth}{@{}>{\bfseries}c>{\raggedright\arraybackslash}p{3.3cm}>{\raggedright\arraybackslash}p{3.0cm}>{\raggedright\arraybackslash}X@{}}
\toprule
\# & Pergunta (resumo) & Opções & Por que e campo \\
\midrule
1 & Você está em perigo agora ou teme pela sua vida? & Sim; Não; Prefiro não responder & Segurança primeiro (OMS). \code{perigo_imediato} \\
2 & O atendimento é para você ou para outra pessoa? & Para mim; Outra pessoa & A Ouvidoria e o 180 recebem relatos de terceiros. \code{para_quem} \\
3 & Qual o assunto principal? & Física; Sexual; Psicológica; Patrimonial; Moral; On-line; Perseguição; Assédio ou importunação; Demora em processo; Outro ou não sei & Taxonomia da Cartilha (p.~6--13) e da Lei 11.340. \code{assunto} \\
4 & Onde ou com quem aconteceu? & Parceiro(a) ou ex; Familiar; Trabalho; Espaço público; Internet; Instituição; Outro; Prefiro não dizer & Define o âmbito (doméstico, familiar, íntimo: art.~5º) e o encaminhamento. \code{ambiente} \\
5 & Já existe processo ou medida protetiva? Se sim, o número (opcional) & Sim, tenho o número; Sim, sem número; Não; Não sei & Escopo processual da Ouvidoria. \code{processo} (informado, não verificado) \\
6 & Posso continuar por este WhatsApp? É seguro? & Sim; Não é seguro; Prefiro telefone & Segurança digital (NNEDV). \code{contato_seguro} \\
7 & Se quiser, conte em poucas palavras. Pode pular. & Texto livre; Pular & Contar uma vez só. \code{relato} (integral, nunca em log) \\
8 & Posso enviar este resumo para a Ouvidoria da Mulher? & Encaminhar; Refazer; Sair sem enviar & Consentimento explícito. \code{consentimento} \\
\bottomrule
\end{tabularx}
\vskip3pt
{\tiny\textcolor{cinza}{Cada pergunta cabe em uma mensagem curta com botões (até 3) ou lista (até 10, título de até 24 caracteres). Fora dos botões, o texto livre só é aceito nos passos 5 e 7; tocar numa opção sempre escapa do campo livre.}}
\end{frame}

\begin{frame}{Caminhos de segurança e exceções}
\begin{columns}[T]
\begin{column}{0.48\textwidth}
\begin{block}{\small Perigo agora = sim}
\scriptsize Antes de qualquer outra pergunta: \textbf{190} (Polícia Militar) e \textbf{180}. O pedido fica marcado como \code{urgencia_imediata}. A Júlia não promete resposta rápida, pois a Ouvidoria atende 8h--14h; pergunta se a pessoa quer registrar o pedido mesmo assim.
\end{block}
\begin{block}{\small Contato não seguro}
\scriptsize Não encaminha ao chat. Mostra telefone e e-mail da Ouvidoria, o 180 e a dica de apagar a conversa. A pessoa pode voltar quando estiver segura. \alert{Decisão sua:} registrar mesmo assim, sem retorno por aqui? (pág.~\ref{pg:decisoes})
\end{block}
\end{column}
\begin{column}{0.48\textwidth}
\begin{block}{\small Fora do horário (8h--14h, Belém)}
\scriptsize Encaminha e avisa que a equipe retoma no próximo horário de atendimento, sem prometer prazo. \alert{Dias da semana: confirmar com a Ouvidoria.}
\end{block}
\begin{block}{\small Falha na transferência}
\scriptsize Reutiliza o fluxo existente: o relato fica guardado, a pessoa escolhe tentar de novo ou ver telefone e e-mail. Nunca diz \emph{encaminhado} sem confirmação.
\end{block}
\begin{block}{\small Quem não é mulher ou pergunta fora do tema}
\scriptsize Não pergunta gênero. Se a pessoa disser que o assunto não é violência contra a mulher, oferece o menu e a Ouvidoria Judiciária.
\end{block}
\end{column}
\end{columns}
\end{frame}

\begin{frame}{Dados estruturados: o que a triagem entrega}
\begin{columns}[T]
\begin{column}{0.47\textwidth}
{\scriptsize\textbf{Esquema v1} (coluna nova \code{triage} no SQLite, validada por enumerações)}
\begin{tikzpicture}
  \node[fill=nevoa,rounded corners=4pt,inner sep=5pt,text width=0.92\textwidth,font=\fontsize{6}{7.2}\selectfont\ttfamily,align=left]{%
\{ "versao": 1,\\
\ "servico": "ouvidoria\_mulher",\\
\ "perigo\_imediato": "sim|nao|nao\_informado",\\
\ "para\_quem": "propria|terceiro",\\
\ "assunto": "fisica|sexual|psicologica|\\
\ \ patrimonial|moral|online|perseguicao|\\
\ \ assedio\_importunacao|morosidade|outro",\\
\ "ambiente": "parceiro\_ex|familiar|trabalho|\\
\ \ espaco\_publico|internet|instituicao|outro",\\
\ "processo": \{"existe": "sim|nao|nao\_sei",\\
\ \ "numero": null, "verificado": false\},\\
\ "contato\_seguro": "sim|nao|telefone",\\
\ "relato": null,\\
\ "consentimento": true \}};
\end{tikzpicture}
\end{column}
\begin{column}{0.50\textwidth}
{\scriptsize\textbf{Nota privada no Chatwoot} (mesma ideia da nota de diligência atual)}
\begin{tikzpicture}
  \node[draw=rosa,fill=rosaclaro,rounded corners=4pt,inner sep=5pt,text width=0.92\textwidth,font=\fontsize{6.5}{8}\selectfont,align=left]{%
\textbf{ATENDIMENTO VIRTUAL: OUVIDORIA DA MULHER}\\
Urgência: \textbf{perigo imediato informado}\\
Para quem: própria \textperiodcentered\ Assunto: violência psicológica\\
Ambiente: parceiro(a) ou ex\\
Processo/medida: sim, nº informado (não verificado)\\
Contato seguro por este WhatsApp: sim\\
Data e horário: (Belém) \textperiodcentered\ Protocolo: TKT-\ldots\\[2pt]
Relato (texto integral): \ldots};
\end{tikzpicture}
\vskip4pt
{\scriptsize \textbf{Rótulos:} \code{ouvidoria_mulher}, \code{urgencia_imediata}, \code{tem_processo}. \textbf{Atributos da conversa} (só enumerações, sem relato): servico, assunto, ambiente, urgencia, para\_quem. O relato e o nº do processo ficam \alert{só na nota e no banco}, nunca em log ou auditoria.}
\end{column}
\end{columns}
\end{frame}

\begin{frame}{O que a triagem não faz}
\begin{columns}[T]
\begin{column}{0.48\textwidth}
\textbf{Não coleta}
\begin{itemize}\small\setlength\itemsep{3pt}
  \item CPF, RG, endereço, nome ou dados do agressor.
  \item Detalhes da violência, além do relato livre opcional.
  \item Gênero ou perfil da pessoa.
  \item Fotos, áudios ou documentos (fora do piloto).
\end{itemize}
\end{column}
\begin{column}{0.48\textwidth}
\textbf{Não faz}
\begin{itemize}\small\setlength\itemsep{3pt}
  \item Aplicar o FONAR, pontuar risco ou classificar \emph{gravidade}.
  \item Dizer que um crime \emph{ocorreu} ou aconselhar processar.
  \item Prometer prazo, resultado, prioridade ou sigilo absoluto.
  \item Confirmar que o processo existe: o acervo do piloto é de demonstração (TJAP).
  \item Insistir: perguntas puladas ficam \emph{não informado}.
\end{itemize}
\end{column}
\end{columns}
\vskip6pt
\begin{tikzpicture}
  \node[fill=areia,rounded corners=5pt,inner sep=6pt,text width=0.93\textwidth,font=\small]{\textbf{Por quê.} O FONAR é do profissional: a própria vítima só o preenche na ausência de profissional (Res.~Conj.~5/2020). Um escore automático teria falsa precisão. A triagem apenas \textbf{sinaliza e organiza}, e quem avalia é a equipe.};
\end{tikzpicture}
\end{frame}

\begin{frame}{Exemplo de conversa (sintético)}
\scriptsize
\begin{columns}[T]
\begin{column}{0.49\textwidth}
\begin{tikzpicture}
  \node[fill=nevoa,rounded corners=6pt,inner sep=5pt,text width=0.94\textwidth,align=left]{\textbf{Júlia:} Aqui você pode falar com calma. Se estiver em perigo agora, ligue \textbf{190}. O \textbf{180} atende 24 h. Quer orientações ou falar com a Ouvidoria da Mulher?};
\end{tikzpicture}
\vskip3pt
\begin{tikzpicture}
  \node[fill=rosaclaro,rounded corners=6pt,inner sep=5pt,text width=0.70\textwidth,align=left,xshift=1.2cm]{\textbf{Cidadã:} Falar com a Ouvidoria};
\end{tikzpicture}
\vskip3pt
\begin{tikzpicture}
  \node[fill=nevoa,rounded corners=6pt,inner sep=5pt,text width=0.94\textwidth,align=left]{\textbf{Júlia:} Você está em perigo agora ou teme pela sua vida? \textcolor{cinza}{[Sim] [Não] [Prefiro não responder]}};
\end{tikzpicture}
\vskip3pt
\begin{tikzpicture}
  \node[fill=rosaclaro,rounded corners=6pt,inner sep=5pt,text width=0.30\textwidth,align=left,xshift=3.4cm]{\textbf{Cidadã:} Não};
\end{tikzpicture}
\end{column}
\begin{column}{0.49\textwidth}
\begin{tikzpicture}
  \node[fill=nevoa,rounded corners=6pt,inner sep=5pt,text width=0.94\textwidth,align=left]{\textbf{Júlia:} Obrigada por responder. Qual é o assunto principal? \textcolor{cinza}{[lista de 10]}\\ \emph{(passos 3 a 6 com toques; passo 7 opcional)}};
\end{tikzpicture}
\vskip3pt
\begin{tikzpicture}
  \node[fill=nevoa,rounded corners=6pt,inner sep=5pt,text width=0.94\textwidth,align=left]{\textbf{Júlia:} Posso enviar este resumo (assunto, ambiente, processo, contato) para a Ouvidoria da Mulher? Eles atendem das 8h às 14h (Belém). \textcolor{cinza}{[Encaminhar] [Refazer] [Sair]}};
\end{tikzpicture}
\vskip3pt
\begin{tikzpicture}
  \node[fill=nevoa,rounded corners=6pt,inner sep=5pt,text width=0.94\textwidth,align=left]{\textbf{Júlia} \emph{(só após a confirmação do Chatwoot)}: Sua conversa foi encaminhada para a Ouvidoria da Mulher. Quando a equipe retomar o atendimento, continuará por aqui. Protocolo: TKT-\ldots};
\end{tikzpicture}
\end{column}
\end{columns}
\vskip2pt
{\tiny\textcolor{cinza}{Redação ilustrativa: os textos finais precisam da revisão da Ouvidoria.}}
\end{frame}

% ==========================================================================
% D. CHATWOOT
% ==========================================================================
\begin{frame}{Novo time ou novo canal?}\label{pg:time}
\scriptsize
\begin{tabularx}{\linewidth}{@{}>{\bfseries}p{3.6cm}>{\raggedright\arraybackslash}X>{\raggedright\arraybackslash}X@{}}
\toprule
Critério & \textcolor{rosa}{Novo time na caixa atual} (recomendado) & Nova caixa de entrada (canal) \\
\midrule
O que é no Chatwoot & Quem responde: agentes agrupados, com atribuição da conversa & Por onde chega: um número ou canal de WhatsApp \\
Roteamento & \code{assign_team} já usado pelo código; basta o nome do time & Exige novo número do WhatsApp Business e novo vínculo do Agent Bot \\
Visibilidade restrita só à Ouvidoria & \alert{Não garante}: em geral, agentes veem as conversas das caixas de que fazem parte (confirmar na versão instalada) & Sim: só os agentes da caixa \\
Esforço e risco & Baixo, reversível, sem mexer no que funciona & Maior: número, templates, vínculo do bot (exige seu pedido explícito) \\
Adequação ao piloto e à demonstração & Suficiente & Desnecessária para a demonstração \\
\bottomrule
\end{tabularx}
\vskip5pt
\begin{tikzpicture}
  \node[fill=areia,rounded corners=5pt,inner sep=6pt,text width=0.93\textwidth,font=\small]{\textbf{Recomendação.} Começar com o \textbf{time} \code{ouvidoria da mulher} (o Chatwoot grava nomes de time em minúsculas). Migrar para uma \textbf{caixa própria} só se a Ouvidoria ou a DPO exigirem \textbf{isolamento de acesso} ou \textbf{número dedicado}. É uma decisão de privacidade, não técnica.};
\end{tikzpicture}
\end{frame}

\begin{frame}{O que fazer no Chatwoot}
\scriptsize
\begin{tabularx}{\linewidth}{@{}>{\bfseries}c>{\raggedright\arraybackslash}p{4.3cm}>{\raggedright\arraybackslash}X@{}}
\toprule
\# & Ação & Detalhe \\
\midrule
1 & Criar o time \code{ouvidoria da mulher} & Configurações $\to$ Equipes. Ou pelo script existente, em modo plano. \\
2 & Adicionar só as pessoas da Ouvidoria & Quem tem acesso a conversas sensíveis é decisão de privacidade; apenas agentes necessários. \\
3 & Criar 3 rótulos & \code{ouvidoria_mulher}, \code{urgencia_imediata}, \code{tem_processo} (Configurações $\to$ Rótulos). \\
4 & Criar atributos personalizados de conversa & servico, assunto, ambiente, urgencia, para\_quem (só enumerações; sem relato). \\
5 & Respostas prontas para a equipe & Acolhimento, confirmação de contato seguro, pedido de nº do processo, dados de 180 e 190. \alert{Quem redige é a Ouvidoria.} \\
6 & Conferir horário e avisos da caixa & Não alterar o horário da caixa inteira: a regra de 8h--14h fica no módulo, só para este serviço. \\
7 & Revisar o acesso a esta caixa & Se a versão instalada tiver funções personalizadas, avaliar; senão, considerar a caixa própria (pág.~anterior). \\
8 & \alert{Não fazer} & Não ligar nem desligar o Agent Bot, não criar caixa e não mexer em \code{secretaria} e \code{gabinete} sem pedido seu. \\
\bottomrule
\end{tabularx}
\vskip3pt
{\tiny\textcolor{cinza}{O script de preparação (\code{chatwoot_setup}) hoje cria times e rótulos e roda em \emph{plano} por padrão; falta suporte a atributos personalizados (item de implementação).}}
\end{frame}

% ==========================================================================
% E. IMPLEMENTAÇÃO
% ==========================================================================
\begin{frame}{Componentes a alterar}
\fontsize{7.2}{8.6}\selectfont
\begin{tabularx}{\linewidth}{@{}>{\raggedright\arraybackslash}p{5.4cm}>{\raggedright\arraybackslash}X@{}}
\toprule
Arquivo & Mudança \\
\midrule
\code{backend/app/domain/triage_mulher.py} \textbf{(novo)} & Enumerações, dataclass da triagem, validação do esquema e geração do resumo; domínio puro, sem I/O. \\
\code{application/chat/flow.py} & Tipos de ação \code{MULHER_*}, opções (títulos dentro dos limites), textos, linha nova no menu principal sob a chave de configuração, reconhecimento por título. \\
\code{application/chat/handler.py} & Métodos \code{_mulher_*}: passos, \code{awaiting="mulher:PASSO"}, perigo, contato inseguro, resumo, nota de contexto, escolha do time. \\
\code{infrastructure/sqlite/operational_store.py} & Coluna \code{triage TEXT} (migração automática \code{ALTER TABLE}), campo no estado, limpeza ao encerrar. \\
\code{app/config.py} e \code{.env.example} & \code{ouvidoria_mulher_enabled=false}, \code{handoff_team_mulher}, limite do relato, horário e contatos. \\
\code{backend/config/mulher.yaml} \textbf{(novo)} & Contatos com \emph{fonte e data de verificação} (mesmo padrão do \code{units.yaml}); nada fabricado. \\
\code{infrastructure/chatwoot/client.py} & Métodos para atributos personalizados da conversa e suas definições. \\
\code{interfaces/chatwoot_setup.py} & Time, 3 rótulos e atributos, em plano por padrão. \\
\code{application/answering/*} e \code{sources.yaml} & Cartilha como fonte institucional \code{pending_review}; busca da \emph{Outra dúvida} restrita a ela; sem aconselhamento. \\
\code{tests/unit}, \code{tests/integration} & Triagem, caminhos de segurança, menu, privacidade, regressão com a chave desligada. \\
\code{docs/OPERACAO.md} & Procedimento de ativação, desativação e revisão. \\
\bottomrule
\end{tabularx}
\end{frame}

\begin{frame}{Plano de implementação}
\scriptsize
\begin{tabularx}{\linewidth}{@{}>{\bfseries}p{1.7cm}>{\raggedright\arraybackslash}X>{\raggedright\arraybackslash}p{2.6cm}@{}}
\toprule
Fase & Tarefas & Depende de \\
\midrule
0. Decisões & D1--D9 (pág.~\ref{pg:decisoes}); autorização para começar & você \\
1. Conteúdo & Obter o PDF da Cartilha e ingerir com OCR como \code{pending_review}; redigir textos fixos; Ouvidoria revisa; confirmar 180, 190 e horário em fonte oficial & Fase 0 \\
2. Chatwoot & Plano (\emph{dry-run}) do script; aplicar time, rótulos e atributos \textbf{com sua aprovação}; membros e respostas prontas & Fase 0 \\
3. Código & Especificação e tarefas (\code{.specs/features/ouvidoria-mulher}); domínio $\to$ armazenamento $\to$ fluxo $\to$ handler $\to$ Chatwoot $\to$ configuração; tarefas atômicas, cada uma com teste & Fases 1 e 2 (textos) \\
4. Homologação & Testes offline; teste em caixa ou número de teste; ensaio com a Ouvidoria; revisão de privacidade (DPO) & Fase 3 \\
5. Entrega & Deploy só de código, chave \textbf{desligada}; ligar na VM somente com seu pedido; Agent Bot só com pedido explícito & Fase 4 \\
\bottomrule
\end{tabularx}
\vskip3pt
{\tiny\textcolor{cinza}{Sem estimativa de prazo: depende das respostas da Ouvidoria (Fase 1) e da sua aprovação (Fases 0 e 2). Tamanho relativo: Fase 3 é a maior; as demais são pequenas ou dependem de terceiros.}}
\end{frame}

\begin{frame}{Testes e critérios de aceite (exemplos)}
\begin{columns}[T]
\begin{column}{0.56\textwidth}
\scriptsize\textbf{Critérios (formato EARS)}
\begin{itemize}\setlength\itemsep{2pt}
  \item \textbf{Quando} a pessoa escolher a Ouvidoria, \textbf{o sistema deve} mostrar 190 e 180 na primeira mensagem.
  \item \textbf{Quando} responder \emph{sim} a perigo, \textbf{deve} enviar o cartão de emergência antes de qualquer outra pergunta e marcar urgência.
  \item \textbf{Quando} tocar \emph{Prefiro não responder}, \textbf{deve} gravar \code{nao_informado} e seguir.
  \item \textbf{Quando} confirmar o envio, \textbf{deve} gravar a nota \textbf{antes} de atribuir o time, e só avisar \emph{encaminhado} após a confirmação do Chatwoot.
  \item \textbf{Se} o contato não for seguro, \textbf{não deve} encaminhar sem novo consentimento.
  \item \textbf{O sistema não deve} registrar relato nem número de processo em log ou auditoria.
  \item \textbf{Com a chave desligada}, o menu deve ser idêntico ao atual.
\end{itemize}
\end{column}
\begin{column}{0.42\textwidth}
\scriptsize\textbf{Cenários sintéticos de desenvolvimento}
\begin{itemize}\setlength\itemsep{2pt}
  \item Violência doméstica, sem perigo, com processo
  \item Perigo agora
  \item Relato de terceira pessoa
  \item Assédio no trabalho
  \item Violência on-line
  \item Demora em processo
  \item Contato inseguro
  \item Pula todas as perguntas
  \item Interrompe com pergunta de outro tema
  \item Fora do horário; falha de transferência
\end{itemize}
{\tiny\textcolor{cinza}{São testes de desenvolvimento. A validação real é o ensaio com a Ouvidoria; os resultados serão informados separadamente, sem misturar os dois.}}
\end{column}
\end{columns}
\end{frame}

\begin{frame}{Riscos e mitigação}
\scriptsize
\begin{tabularx}{\linewidth}{@{}>{\raggedright\arraybackslash}p{4.3cm}>{\raggedright\arraybackslash}X@{}}
\toprule
Risco & Mitigação \\
\midrule
Pessoa real em perigo usa o número da demonstração & Piloto sinalizado; 190 e 180 em toda tela; pergunta de perigo primeiro; a Ouvidoria define quem monitora e quando. \\
Agressor vê o celular ou as notificações & Pergunta de contato seguro; saída rápida; aviso de apagar a conversa; mensagens sem detalhes sensíveis. \\
Triagem vira obstáculo & \emph{Falar com atendente} a qualquer momento; todas as perguntas têm saída. \\
Falso positivo no atalho por texto livre & Só oferece o serviço; não força. \\
Dado sensível exposto a muitos agentes & Time restrito; avaliar caixa própria; relato fora de log; política de retenção definida com a DPO. \\
Conteúdo errado ou desatualizado & Textos aprovados pela Ouvidoria; contatos com fonte e data; Cartilha em \code{pending_review} até um humano aprovar. \\
Horário 8h--14h cria expectativa falsa & Mensagem clara de horário; sem prazo prometido. \\
\bottomrule
\end{tabularx}
\end{frame}

\begin{frame}{Decisões de que preciso}\label{pg:decisoes}
\small
\begin{enumerate}\setlength\itemsep{3pt}
  \item \textbf{D1} Aprova o \textbf{time} \code{ouvidoria da mulher} (em vez de nova caixa)? Quem são os membros?
  \item \textbf{D2} O número de WhatsApp atual é aceitável para o assunto, ou a Ouvidoria/DPO exige \textbf{caixa e número próprios}?
  \item \textbf{D3} \textbf{Contato não seguro:} apenas orientar, ou permitir registrar o pedido sem retorno pelo WhatsApp?
  \item \textbf{D4} \textbf{Horário:} confirmar 8h--14h e os dias da semana; mensagem fora do horário.
  \item \textbf{D5} Aceita \textbf{relatos de terceiros} (testemunha, familiar)? Proposta: sim.
  \item \textbf{D6} \textbf{Quem aprova} os textos das orientações e a ingestão da Cartilha? Pode enviar o PDF para o repositório?
  \item \textbf{D7} \textbf{Retenção} do relato no banco do piloto (tempo e eliminação), com a DPO.
  \item \textbf{D8} \textbf{Confirmar contatos}: WhatsApp do 180, 190 e dados da Ouvidoria, em fonte oficial atual.
  \item \textbf{D9} Posso criar a especificação e começar a implementar atrás da chave desligada? E o menu: posição 4, como proposto?
\end{enumerate}
\end{frame}

\begin{frame}{Referências}
\tiny
\begin{itemize}\setlength\itemsep{2pt}
  \item Cartilha da Ouvidoria da Mulher (TJPA), PDF fornecido: págs.~4--21.
  \item TJPA, \emph{Carta de Serviços}, ficha Ouvidoria da Mulher: \url{https://www.tjpa.jus.br/CMSPortal/VisualizarArquivo?idArquivo=1292557}
  \item TJPA, página da Ouvidoria da Mulher (indicada na Carta): \url{https://www.tjpa.jus.br/PortalExterno/institucional/Ouvidoria-da-Mulher/982289-apresentacao.xhtml}
  \item CNJ/CNMP, Resolução Conjunta 5/2020 (FONAR) e Lei 14.149/2021; guia: \url{https://dtic.tjpr.jus.br/documents/12054912/56756885/Formulário+Nacional+de+Avaliação+de+Risco.pdf/49c39f8c-ee93-054f-5dbe-c854178a0f48}
  \item OMS, \emph{Health care for women subjected to intimate partner violence or sexual violence: a clinical handbook}: \url{https://www.who.int/publications/i/item/9789241517102} (currículo associado)
  \item Campbell, Webster e Glass (2009), \emph{The Danger Assessment}, \emph{J. Interpersonal Violence}; Storey (2014).
  \item NNEDV, \emph{Technology Abuse: Experiences of Survivors and Victim Service Agencies}: \url{https://nnedv.org/latest_update/technology-abuse-experiences-of-survivors-and-victim-service-agencies/}
  \item Lei 14.550/2023 (Cartilha p.~16): \url{https://www.lex.com.br/lei-no-14-550-de-19-de-abril-de-2023}
  \item CNJ, Resolução 254/2018 (política judiciária de enfrentamento): \url{https://atos.cnj.jus.br/atos/detalhar/2669}
  \item Ligue 180 e 190 (fonte secundária; \alert{confirmar em gov.br}): \url{https://fpabramo.org.br/presidente-lula-sanciona-tres-leis-para-proteger-a-vida-das-brasileiras/}
  \item Código do projeto: \code{backend/app/application/chat/{flow,handler}.py}, \code{domain/handoff.py}, \code{interfaces/chatwoot_setup.py}.
\end{itemize}
\end{frame}

\end{document}
